\subsection{SEC reporting: the lead profile survives the design; the atypicality gradient is composition}\label{sec:bluesea}

The outcome here is the firm-level one that the representation hazard
of Section~\ref{sec:hazard} turns on: whether an owner ever appears in the Securities and
Exchange Commission's financial reporting records. That is broader than
exchange-listed: roughly half the matched firms carry a ticker and the
rest report for other reasons. Figure~\ref{fig:successcurves} draws the
raw rate against each axis in twenty bins. In signed lead the profile is a
U whose tails both sit above the base rate and whose lagging tail is the
higher; in lead magnitude and in atypicality the raw rate rises.

\begin{figure}[h]\centering
\includegraphics[width=\textwidth]{results/fig_success_ventiles.png}
\caption{Share of registered debut owners (first registered filing
1995--2018, $n = 2{,}066{,}600$) ever appearing in SEC EDGAR, in twenty
equal-count bins of debut lead $L$, lead magnitude $|L|$, and atypicality
$A$, cut on the pooled distribution. Bands are 95\% binomial intervals;
the dashed line is the 0.65\% base rate. These are raw rates carrying
class, year and length composition; Table~\ref{tab:edgar} gives the
design-based estimates, under which the atypicality gradient is absorbed
and the signed-lead U remains.}
\label{fig:successcurves}
\end{figure}

The raw atypicality gradient does not survive the design-based estimates
(Table~\ref{tab:edgar}). Within class and filing year, holding description
length fixed, top-quintile atypicality raises the probability of SEC
reporting by $+0.000$pp on a 0.481\% base ($t = 0.0$), and by $-0.003$pp
on the past level alone. The quintile profile is not even monotone: the
three middle quintiles sit \emph{below} the bottom one and the top merely
returns to it. The 2.7-fold gradient in the raw bins is composition:
firms that go on to report cluster in particular classes, years, and
multi-class filings, and class$\times$year fixed effects absorb it. Lead
costs, at equal atypicality: the signed component enters negatively
($-0.047$pp per standard deviation, $t = -7.8$) while the unsigned
component enters positively ($+0.067$pp per standard deviation,
$t = 7.7$). The U in the raw bins is exactly this pair: lagging filings
show the magnitude premium with no lead penalty, leading filings show it
net of that penalty, and the middle bins show neither. Once the levels
are controlled, the recovery of the far-leading tail is much reduced: the
quintile profile falls through Q4 and recovers only slightly at Q5.

\begin{table}[h]\centering\small
\caption{Debut listing and vocabulary position: design-based estimates.
LPM of P(owner ever in SEC EDGAR) on 1{,}556{,}968 registered debut
filings 1995--2018, one observation per owner, class$\times$filing-year
fixed effects, heteroskedasticity-robust SEs. Quintiles cut within
class$\times$year. Base rate 0.481\%.}
\label{tab:edgar}
\begin{tabular}{lrr}
\toprule
 & Coefficient (pp) & SE \\
\midrule
\multicolumn{3}{l}{\emph{Atypicality quintiles, $A = \tfrac{1}{2}(K^{-}+K^{+})$ (Q1 omitted)}} \\
\quad Q5, FE only & $+0.017$ & $0.018$ \\
\quad Q5, FE + log length & $+0.000$ & $0.018$ \\
\quad Q5, FE + log length, on $K^{-}$ alone & $-0.003$ & $0.018$ \\
\midrule
\multicolumn{3}{l}{\emph{Signed lead quintiles, FE + log length + KL levels}} \\
\quad Q3 & $-0.133$ & $0.020$ \\
\quad Q4 & $-0.187$ & $0.021$ \\
\quad Q5 & $-0.158$ & $0.024$ \\
\midrule
\multicolumn{3}{l}{\emph{Decomposition, FE + log length}} \\
\quad lead magnitude $|L|$ (pp per sd) & $+0.067$ & $0.009$ \\
\quad signed lead $L$ (pp per sd) & $-0.047$ & $0.006$ \\
\quad log description length & $+0.056$ & $0.004$ \\
\bottomrule
\end{tabular}
\\[0.3em]\footnotesize Standard errors in place of significance stars; see
the note to Table~\ref{tab:gate_lpm}.
\end{table}

Owners are matched to SEC registrants by normalized exact name, over a
candidate universe of all 5.67 million owners in the corpus; names
resolving to more than one registrant are dropped rather than assigned,
and one match is kept per owner string. That yields 19{,}889 matched
owners covering 7{,}588 registrants, and a base rate of 0.481\% among
registered debut filers. Validation on 22 companies that indisputably went
public --- across software, pharmaceuticals, food, apparel, aerospace and
autos --- matches all 22. A random sample of 25 matches carrying at least
three filings was inspected by hand and none was wrong, which is a check
on gross failure modes rather than a precision estimate.

The linkage is nonetheless a name match, and that bounds the atypicality
reading in particular. Firms with distinctive names are easier to match
than firms with generic ones, and distinctiveness of a name plausibly
travels with distinctiveness of the goods/services text; nothing in this
design separates the two. The lead coefficient is unchanged under
alternative linkages; the atypicality coefficient is not, so only the
lead result is robust to the name-match design.

The U is also not atypicality working through lead's construction (the
mechanical coupling stated in Section~\ref{sec:three_scorings}): on this
sample lead magnitude and atypicality correlate at $+0.25$, signed lead at
only $-0.10$, and cutting lead fifths within fifths of atypicality keeps
the structure: the most lagging fifth out-reports the most leading in four
of the five atypicality strata (by $0.08$ to $0.39$pp on sub-1\% bases),
and the ordering reverses only in the most atypical fifth, where the most
leading debuts report most. The lagging tail's premium is not a ceiling
artifact; the leading tail's partial recovery in the raw U is where the
atypicality composition lives, which is what the level controls in
Table~\ref{tab:edgar} absorb.

``SEC reporting'' is not ``went public'': splitting matched owners on
whether their SEC filer number carries a ticker gives $10{,}556$ listed
against $9{,}333$ reporting without one, and both halves show the same U
in lead (Section~\ref{sec:robust}), so the result is not an artifact of
counting non-operating registrants as successes.

